\documentclass[11pt,letterpaper]{article}

\usepackage[top=1in, bottom=1in, left=1in, right=1in, headheight=14pt]{geometry}
\usepackage{fontspec}
\setmainfont{DejaVu Serif}
\setsansfont{DejaVu Sans}
\setmonofont{DejaVu Sans Mono}

\usepackage{xcolor}
\usepackage{titlesec}
\usepackage{fancyhdr}
\usepackage{enumitem}
\usepackage{hyperref}
\usepackage{booktabs}
\usepackage{tabularx}
\usepackage{longtable}
\usepackage{colortbl}
\usepackage{multirow}
\usepackage{array}
\usepackage{parskip}
\usepackage{tcolorbox}
\usepackage{graphicx}
\usepackage{lastpage}

% ── Color Scheme: Studio Teal / Electric Blue / Warm Gold ──
\definecolor{primary}{HTML}{004D5E}
\definecolor{secondary}{HTML}{0277BD}
\definecolor{tertiary}{HTML}{8B6914}
\definecolor{accent}{HTML}{00838F}
\definecolor{lightprimary}{HTML}{E0F7FA}
\definecolor{lightsecondary}{HTML}{E3F2FD}
\definecolor{lighttertiary}{HTML}{FFF8E1}
\definecolor{rowgray}{HTML}{F5F5F5}
\definecolor{bordergray}{HTML}{BDBDBD}
\definecolor{subtlegray}{HTML}{757575}
\definecolor{darktext}{HTML}{212121}

% ── Section Formatting ──
\titleformat{\section}
  {\Large\bfseries\sffamily\color{primary}}
  {\thesection}{1em}{}
  [\vspace{-0.5em}\textcolor{bordergray}{\rule{\textwidth}{0.4pt}}]

\titleformat{\subsection}
  {\large\bfseries\sffamily\color{secondary}}
  {\thesubsection}{1em}{}

\titleformat{\subsubsection}
  {\normalsize\bfseries\sffamily\color{tertiary}}
  {\thesubsubsection}{1em}{}

\titlespacing*{\section}{0pt}{1.5em}{0.8em}
\titlespacing*{\subsection}{0pt}{1.2em}{0.5em}
\titlespacing*{\subsubsection}{0pt}{0.8em}{0.3em}

% ── Custom Environments ──
\newcommand{\stageheader}[2]{%
  \noindent\colorbox{#2}{\makebox[\textwidth][l]{%
    \textcolor{white}{\bfseries\sffamily\hspace{6pt}#1\hspace{6pt}}}}%
  \vspace{0.5em}\par%
}

\newtcolorbox{asciibox}{
  colback=rowgray, colframe=bordergray,
  arc=1pt, boxrule=0.5pt,
  left=6pt, right=6pt, top=4pt, bottom=4pt,
  fontupper=\small\ttfamily,
  breakable
}

\newtcolorbox{quotebox}{
  colback=lightprimary, colframe=primary,
  arc=2pt, boxrule=0.5pt,
  left=12pt, right=12pt, top=8pt, bottom=8pt,
  breakable
}

\newcommand{\metafield}[2]{\textbf{\sffamily #1} #2\par}

% ── Table Helpers ──
\newcolumntype{L}[1]{>{\raggedright\arraybackslash}p{#1}}
\newcolumntype{C}[1]{>{\centering\arraybackslash}p{#1}}
\newcommand{\tableheader}[1]{\textbf{\sffamily\textcolor{white}{#1}}}

% ── Headers / Footers ──
\pagestyle{fancy}
\fancyhf{}
\fancyhead[L]{\small\sffamily\textcolor{subtlegray}{Ableton Live Research Mission}}
\fancyhead[R]{\small\sffamily\textcolor{subtlegray}{GSD Mission Package}}
\fancyfoot[C]{\small\sffamily\textcolor{subtlegray}{Page \thepage\ of \pageref{LastPage}}}
\renewcommand{\headrulewidth}{0.4pt}
\renewcommand{\footrulewidth}{0pt}

% ── Hyperlinks ──
\hypersetup{
  colorlinks=true,
  linkcolor=secondary,
  urlcolor=secondary,
  pdfauthor={GSD Ecosystem / Tibsfox},
  pdftitle={Ableton Live -- Architecture, Extensibility \& Ecosystem -- Mission Package},
  pdfsubject={Vision to Mission Pipeline},
}

% =========================================================
\begin{document}
% =========================================================

% ── Title Page ──
\thispagestyle{empty}
\begin{center}
\vspace*{3cm}

{\small\sffamily\textcolor{primary}{\textbf{GSD RESEARCH MISSION PACKAGE}}}

\vspace{0.3cm}
\textcolor{bordergray}{\rule{8cm}{0.4pt}}

\vspace{1.5cm}
{\fontsize{28}{34}\selectfont\bfseries\sffamily\textcolor{primary}{Ableton Live\\[0.3em]Architecture, Extensibility \& Ecosystem}}

\vspace{0.8cm}
{\Large\sffamily\textcolor{secondary}{Digital Audio Workstation --- Deep Research Study}}

\vspace{0.5cm}
{\large\sffamily\itshape\textcolor{tertiary}{From Clip Paradigm to Live Object Model}}

\vspace{3cm}
{\sffamily\textcolor{accent}{Vision $\rightarrow$ Research $\rightarrow$ Mission}}\\[0.3em]
{\small\sffamily\textcolor{accent}{Full Pipeline Execution}}

\vspace{3cm}
{\small\sffamily\textcolor{subtlegray}{Date: March 27, 2026  |  Status: Ready for GSD Orchestrator}}\\[0.3em]
{\small\sffamily\textcolor{subtlegray}{Activation Profile: Squadron (12 roles)  |  Estimated: 8--12 context windows across 3 sessions}}
\end{center}
\newpage

% ── Table of Contents ──
\tableofcontents
\newpage

% =========================================================
% STAGE 1 — VISION DOCUMENT
% =========================================================
\stageheader{STAGE 1 --- VISION DOCUMENT}{primary}

\section{Ableton Live --- Vision Guide}

\metafield{Date:}{March 27, 2026}
\metafield{Status:}{Research Complete / Ready for Mission Execution}
\metafield{Depends On:}{Deep Audio Analyzer Mission (prior GSD artifact)}
\metafield{Informs:}{College of Knowledge Audio Module, GSD-OS Audio Integration}
\metafield{Version:}{v1.0}

\subsection{Vision}

Ableton Live is not a conventional digital audio workstation. It is, at its core, a philosophical
statement about the relationship between improvisation and composition---the space between
a free-flowing jam and a committed arrangement. The software's founders, Gerhard Behles and
Robert Henke, embedded that tension directly into the interface: two views, Session and
Arrangement, sharing the same tracks but offering radically different relationships with time.

Where nearly every other DAW imposes a linear timeline as the primary mode of
working---treating music-making as a left-to-right accumulation of fixed decisions---Live's
Session View holds open a non-linear field of possibility. Clips launch on demand. Scenes
trigger in combinations. The music breathes without committing. Only when the producer is
ready does the Arrangement View crystallize that improvisation into a linear form.

This duality maps directly onto the ``spaces between'' that animate the GSD ecosystem. The
Session View is the world of pure potential---the gap before commitment, where combinations
form and dissolve. The Arrangement View is where that potential collapses into execution. The
Tab key that switches between them is among the most philosophically loaded keystrokes in
music software.

Live 12 (released early 2024) extends this philosophy into new territories: generative MIDI
transformations, alternative tuning systems, MPE-native instruments like Meld and an updated
Granulator III, and a sophisticated extensibility layer---Max for Live, Python MIDI Remote
Scripts, and the AbletonOSC protocol---that allows the DAW to be reshaped from outside as
readily as from within. Understanding this extensibility layer is the key to integrating Live
deeply with the GSD ecosystem's audio analysis and creative automation capabilities.

\subsection{Problem Statement}

\begin{enumerate}
  \item \textbf{Dual-View Cognitive Load.} The Session/Arrangement paradigm is Ableton's
  greatest strength and its steepest learning curve. New users consistently struggle to
  understand why clips launched in Session do not affect the Arrangement timeline, and
  vice versa. The ``Back to Arrangement'' button---a small control with enormous behavioral
  consequences---trips up experienced users for months. A structured model of the
  two-view interaction resolves this and unlocks the paradigm's creative power.

  \item \textbf{Synthesis Depth vs.\ Accessibility.} Live 12 ships with 20 instruments spanning
  physical modeling (Analog, Collision, Tension), FM synthesis (Operator), wavetable
  (Wavetable), bi-timbral macro oscillators (Meld), subtractive MPE (Drift), granular
  (Granulator III), and sampling (Sampler, Simpler). The breadth is extraordinary but
  opaque. There is no cohesive synthesis taxonomy guiding a producer from one architecture
  to the next, nor a map of which instrument fits which sound-design intent.

  \item \textbf{Extensibility Layer Fragmentation.} Three parallel extension mechanisms
  co-exist without clear hierarchy: Max for Live (visual patching with full API access),
  Python MIDI Remote Scripts (compiled bytecode interface), and AbletonOSC (network-based
  OSC control). Each accesses the Live Object Model differently. A unified map of the
  extensibility stack---and guidance on when to use which path---is absent from official
  documentation.

  \item \textbf{Hardware Ecosystem Complexity.} Push 3 ships in two configurations---standalone
  and controller---powered by an Intel NUC compute element and running a subset of Live 12
  Suite. Move (October 2024) adds a compact groovebox tier. Ableton Link connects these and
  third-party apps over local networks. Understanding the hardware/software boundary, its
  current limitations (no Arrangement View in Push standalone as of 2025), and the
  community workarounds is essential for workflow design.

  \item \textbf{Education Gap at Scale.} Ableton's Certified Trainer program spans 50 countries
  and 38 languages, with more than 360 trainers. Yet formal documentation of Live's deeper
  architectural decisions---why Session View works the way it does, why the clip paradigm
  was chosen, how the Live Object Model is structured---is scattered across forum posts,
  decompiled Python bytecode, and third-party tutorials. A cohesive reference synthesizing
  architecture, instruments, extensibility, and educational pathways does not yet exist.
\end{enumerate}

\subsection{Core Concept}

\textbf{The Clip Paradigm as Architectural Foundation}

Live's fundamental unit is the clip---a bounded, loopable segment of audio or MIDI. Everything
in Live's architecture flows from this choice. Session View is a two-dimensional field of clips
organized into tracks (vertical) and scenes (horizontal). Arrangement View is the same clips
arranged on a linear timeline. The audio engine treats both identically; only the launching
mechanism differs.

\subsubsection{Dual-View Architecture (ASCII)}

\begin{asciibox}
SESSION VIEW                          ARRANGEMENT VIEW
============                          ================

Track 1 | [Clip A] [Clip B] [Clip C]  Track 1 |----[A]------[B]---[C]---->
Track 2 | [Clip D] [Clip E] [   ]     Track 2 |--------[D]----[E]-------->
Track 3 | [Clip F] [   ]   [Clip G]   Track 3 |--[F]-----------[G]------->
                                                         time --------->

Scene 1 = Launch A+D+F simultaneously
Scene 2 = Launch B+E simultaneously
Scene 3 = Launch C+G simultaneously

[ Tab key ] switches views. Tracks are SHARED.
Session clips and Arrangement clips are MUTUALLY EXCLUSIVE per track.
\end{asciibox}

\subsubsection{Module Architecture}

\begin{asciibox}
  +-----------+   +-----------+   +------------+   +-----------+   +-----------+
  | MODULE 01 |   | MODULE 02 |   | MODULE 03  |   | MODULE 04 |   | MODULE 05 |
  | Core Arch |   | Synthesis |   | Hardware   |   | Extension |   | Education |
  | & Engine  |   | Arsenal   |   | Ecosystem  |   | Layer     |   | Community |
  +-----------+   +-----------+   +------------+   +-----------+   +-----------+
       |               |               |               |               |
       +-------+-------+               +-------+-------+               |
               |                               |                       |
         WAVE 1-A                         WAVE 1-B                 WAVE 1-C
    (Architecture +               (Hardware + Extension)        (Education)
      Instruments)
               \                               /                       /
                +------------- WAVE 2 ---------+---------------------+
                               (Synthesis: Cross-module integration)
\end{asciibox}

\subsection{Research Modules}

\subsubsection{Module 01 --- Core Architecture and Audio Engine}
The fundamental clip paradigm; Session View vs.\ Arrangement View behavioral model;
audio warping engine; MIDI routing; automation; Warp Modes; the ``Back to Arrangement''
state machine; global quantization; Ableton's edition tiers (Intro / Standard / Suite).

\subsubsection{Module 02 --- Instruments and Sound Design Arsenal}
All 20 instruments in Live 12 Suite: physical modeling (Analog, Collision, Electric, Tension),
FM (Operator), wavetable (Wavetable), subtractive MPE (Drift), bi-timbral macro oscillators
(Meld), granular (Granulator III), sampling (Sampler, Simpler), drums (Drum Rack, Drum Synths,
Impulse), bass (Bass, Poli); Live 12 additions (Meld, Roar); MPE capability taxonomy.

\subsubsection{Module 03 --- Hardware Ecosystem}
Push 3 (Controller and Standalone configurations); Intel NUC compute element; MPE pad grid;
built-in audio interface; ADAT I/O; Ableton Move (October 2024 groovebox); Ableton Link
protocol; third-party hardware integration (Akai APC, Novation Launchpad, Push 1/2 legacy).

\subsubsection{Module 04 --- Extensibility Layer}
Max for Live (M4L) and the Live Object Model (LOM); Python MIDI Remote Scripts (bytecode
architecture, Python 3 in Live 11+); AbletonOSC (OSC-based full API access via ideoforms);
ClyphX Pro; Remotify; community API documentation resources (NSUSpray, Julien Bayle);
extension mechanism decision tree.

\subsubsection{Module 05 --- Education, Community and Ecosystem}
Ableton Certified Trainer program (360+ trainers, 50 countries, 38 languages, since 2008);
key training institutions (343 Labs NYC/Berlin, Point Blank, Delta Sound Labs/Case Western);
Ableton's free online learning resources; community platforms; the ``Learning Music'' browser
app; festival culture influence; historical context (Monolake / Robert Henke's departure 2016).

\subsection{Chipset Configuration}

\begin{asciibox}
# ableton-live-mission-chipset.yaml
# GSD Amiga Chipset Assignment

agnus:                        # Orchestration / DMA / Context
  model: claude-opus-4-6
  role: FLIGHT + INTEG
  responsibilities:
    - Mission direction and wave gate decisions
    - Cross-module synthesis (architecture --> instruments --> hardware)
    - Live Object Model deep analysis
    - GSD ecosystem connection analysis

denise:                       # Display / Output Rendering
  model: claude-sonnet-4-6
  role: EXEC-1 + EXEC-2 + CRAFT-ARCH + CRAFT-EXT
  responsibilities:
    - Module 01 document production (Core Architecture)
    - Module 02 instrument taxonomy compilation
    - Module 04 extensibility layer mapping
    - Final document assembly and formatting

paula:                        # Audio / I-O / Anticipatory
  model: claude-haiku-4-5
  role: SCOUT + BUDGET + LOG
  responsibilities:
    - Source gathering and gap-fill web research
    - Token budget tracking
    - Shared schema and template scaffolding
    - Commit messages and audit trail

68000:                        # CPU / Glue Logic
  model: claude-sonnet-4-6
  role: PLAN + VERIFY + SURGEON + CAPCOM
  responsibilities:
    - Wave decomposition and parallel track management
    - Source validation and accuracy checking
    - Research quality degradation detection
    - Human approval at wave boundaries
\end{asciibox}

\subsection{Scope Boundaries}

\subsubsection{In Scope (v1.0)}
\begin{itemize}[leftmargin=1.5em]
  \item Ableton Live 12 (current release as of March 2026)
  \item All five research modules as defined above
  \item Live editions: Intro, Standard, Suite (not Lite)
  \item Push 3 (Controller and Standalone), Ableton Move
  \item Max for Live, Python MIDI Remote Scripts, AbletonOSC
  \item Certified Trainer network overview (not individual trainer profiles)
  \item GSD ecosystem integration opportunities identified per module
  \item Ableton Link protocol and network synchronization
\end{itemize}

\subsubsection{Out of Scope}
\begin{itemize}[leftmargin=1.5em]
  \item Third-party VST/AU plugin ecosystem (Serum, Massive, Omnisphere, etc.)
  \item Competitive DAW comparison (Logic, Pro Tools, FL Studio, Bitwig)
  \item Audio mastering workflows in depth
  \item Ableton's business/financial history
  \item Push 1 and Push 2 (legacy hardware) beyond brief historical mention
  \item Individual artist production technique analysis
  \item Video production or scoring-to-picture workflows
\end{itemize}

\subsection{Success Criteria}

\begin{enumerate}
  \item Session View and Arrangement View behavioral model fully documented, including the
  clip exclusivity rule and ``Back to Arrangement'' state transitions.

  \item All 20 instruments in Live 12 Suite catalogued with synthesis architecture, key
  parameters, and recommended use-cases.

  \item Extensibility mechanism decision tree produced: when to use M4L vs.\ Python Remote
  Scripts vs.\ AbletonOSC for a given integration goal.

  \item Live Object Model hierarchy documented to at least three levels: Song $\to$ Track $\to$
  Clip with key observable properties identified.

  \item Push 3 Standalone limitations (as of 2025/2026) documented alongside community
  workarounds; Ableton Move capabilities enumerated.

  \item Ableton Link protocol behavior documented: tempo negotiation, quantization, start/stop
  independence, network topology.

  \item Education ecosystem mapped: Certified Trainer program scale and structure; top
  institutions; free learning resources available via ableton.com.

  \item GSD integration pathways identified for at least three modules (suggested: Core
  Architecture $\to$ workflow templates; Extensibility $\to$ Python script generation;
  Hardware $\to$ Push/Move as GSD output device).

  \item All instrument and device claims sourced to Ableton's official reference manual
  (Version 12) or verifiable product announcements.

  \item Document is self-contained: a reader with no prior Ableton knowledge can orient
  themselves using Stage 1 alone.

  \item Synthesis taxonomy produced linking instrument architecture to GSD mathematical
  foundations (e.g., Operator's FM synthesis $\to$ complex number modulation from
  \textit{The Space Between}).

  \item Test plan covers all five modules with at least two tests per success criterion
  and all safety-critical source-attribution tests passing.
\end{enumerate}

\subsection{Relationship to Other Documents}

\begin{center}
\rowcolors{2}{rowgray}{white}
\begin{tabularx}{\textwidth}{L{4cm}L{4cm}X}
\toprule
\rowcolor{primary}
\tableheader{Document} & \tableheader{Relationship} & \tableheader{Dependency Direction} \\
\midrule
Deep Audio Analyzer Mission & Prior GSD artifact & This mission extends audio analysis to DAW integration \\
The Space Between (Tibsfox) & Mathematical spine & FM synthesis (Operator) maps to complex number modulation; L-Systems map to generative MIDI \\
GSD-OS Desktop Vision & Platform context & Audio integration layer for GSD-OS uses Live extensibility findings \\
College of Knowledge & Educational layer & Audio module seeded by this research; instrument taxonomy feeds curriculum \\
Skill-Creator (v1.49+) & Execution engine & Mission handed to skill-creator for agent-based execution \\
\bottomrule
\end{tabularx}
\end{center}

\subsection{The Through-Line}

Ableton Live embodies the Amiga Principle at its most musical. A minimal architectural
decision---two views sharing the same tracks---produces a creative topology that has reshaped
how electronic music is made worldwide, contributed to the rise of global festival culture
through the 2000s, and made the laptop a legitimate instrument on stage alongside guitars
and drums. The sophistication is not in the number of features, but in the principled elegance
of a single foundational idea: a clip is a bounded unit of musical potential, and the space
between its launch and its commitment to the timeline is where music lives.

In the GSD framework, this maps precisely to the subagent architecture: each agent operates
on a bounded context (a clip), can be triggered non-linearly (Session View), and eventually
contributes its output to a linear pipeline (Arrangement View). The Tab key between Session
and Arrangement is the same conceptual move as GSD's transition from wave-based exploration
to committed synthesis. Max for Live's extensibility model---where any Live parameter can be
observed, queried, and controlled via the Live Object Model---is a precursor to the kind of
deterministic inter-agent communication that DACP formalizes in the skill-creator ecosystem.
The spaces between the clips are not silence; they are the architecture's most eloquent feature.

% =========================================================
% STAGE 2 — RESEARCH REFERENCE
% =========================================================
\newpage
\stageheader{STAGE 2 --- RESEARCH REFERENCE}{secondary}

\section{Ableton Live --- Research Reference}

\subsection{How to Use This Document}

This reference compiles findings from web research conducted March 27, 2026. All claims are
sourced to Ableton's official documentation (Reference Manual Version 12), official product
announcements, or established audio technology publications. Entertainment media and unsourced
forum opinions are excluded from findings (though community forum context is used to identify
known user pain points).

\textbf{Key source organizations:}
\begin{itemize}[leftmargin=1.5em]
  \item \textbf{Ableton AG} --- Official Reference Manual v12 (ableton.com/en/manual)
  \item \textbf{Ableton Help Center} --- Product support and API documentation
  \item \textbf{CDM Create Digital Music} --- Authoritative electronic music technology press
  \item \textbf{Sound on Sound} --- Professional audio publication
  \item \textbf{Cycling '74} --- Max/MSP developer (Max for Live partner)
  \item \textbf{Wikipedia} --- Ableton Live article for historical facts and edition lists
  \item \textbf{GitHub (ideoforms/AbletonOSC)} --- AbletonOSC open-source project
  \item \textbf{GitHub (gluon/AbletonLive11\_MIDIRemoteScripts)} --- Community API documentation
\end{itemize}

\subsection{Module 01 Reference: Core Architecture and Audio Engine}

\subsubsection{The Session View Paradigm}

Ableton's Session View is designed explicitly for the non-linear creative modes that a
traditional DAW timeline cannot support: live performance, DJing, improvisation, and
score-reactive theatre work. According to the official Reference Manual Version 12, the view
``centers around improvisation, looping, and clip launching'' and allows music to be structured
and restructured without stopping playback.

Each track in Session View contains a vertical column of clip slots. Each clip slot can hold
one audio or MIDI clip. At any moment, at most one clip per track is playing. Scenes---rows
across the Session grid---provide a mechanism for launching multiple clips simultaneously with
a single trigger, enabling full ``song sections'' to fire at once.

The \textbf{global launch quantization} setting is among Session View's most consequential
controls: it determines how long Live waits after a clip trigger before actually starting
playback, ensuring clips launch on a musically sensible boundary (a bar, a beat, etc.)
rather than mid-phrase.

\subsubsection{The Arrangement View Paradigm}

The Arrangement View is a conventional horizontal timeline---analogous to the edit windows in
Pro Tools, Cubase, or Logic---where clips are placed at specific positions in musical time.
Audio and MIDI clips appear as colored blocks on track lanes; the user edits them by dragging,
resizing, and applying automation envelopes.

A critical and often-misunderstood feature governs the interaction between the two views:
\textbf{Session clips and Arrangement clips on any given track are mutually exclusive}---only
one can play at a time. When a Session clip is launched, Live halts that track's Arrangement
playback in favor of the Session clip. The Arrangement only resumes when the user explicitly
clicks the ``Back to Arrangement'' button, which lights up as a reminder that what is being
heard differs from the saved Arrangement state.

\subsubsection{Core Editing Capabilities (Arrangement View)}

\begin{center}
\rowcolors{2}{rowgray}{white}
\begin{tabularx}{\textwidth}{L{4.5cm}X}
\toprule
\rowcolor{primary}
\tableheader{Feature} & \tableheader{Description} \\
\midrule
Clip warping & Audio clips can be time-stretched to any tempo using six Warp Modes \\
Automation & All Live parameters can be automated as breakpoint curves on the timeline \\
Time signature markers & Not quantized; can be placed anywhere; fragmentary bars supported \\
Fade/crossfade & 4\,ms auto-crossfades on adjacent clips; custom fade handles available \\
Selection-based editing & Cut, copy, paste, duplicate operate on timeline selections \\
Waveform zoom & Vertical waveform zoom slider; non-destructive; applies to all audio tracks \\
Locators & Named position markers for navigation during recording and performance \\
\bottomrule
\end{tabularx}
\end{center}

\subsubsection{Edition Tiers}

\begin{center}
\rowcolors{2}{rowgray}{white}
\begin{tabularx}{\textwidth}{L{2.5cm}L{2.5cm}X}
\toprule
\rowcolor{primary}
\tableheader{Edition} & \tableheader{Instruments} & \tableheader{Notable Inclusions / Exclusions} \\
\midrule
Intro & 4 & Impulse, Simpler, Instrument Rack, Drum Rack; limited tracks \\
Standard & 6+ & Adds External Instrument; unlimited tracks; all mixing effects \\
Suite & 20 & All instruments; 58 audio effects; 14 MIDI effects; Max for Live; 71+ GB content \\
Lite & (bundled) & Similar to Intro; bundled with hardware; not sold standalone \\
\bottomrule
\end{tabularx}
\end{center}

\subsection{Module 02 Reference: Instruments and Sound Design Arsenal}

\subsubsection{Synthesis Architecture Taxonomy (Live 12 Suite)}

\begin{center}
\rowcolors{2}{rowgray}{white}
\begin{tabularx}{\textwidth}{L{2.5cm}L{3cm}X}
\toprule
\rowcolor{primary}
\tableheader{Instrument} & \tableheader{Architecture} & \tableheader{Key Characteristics} \\
\midrule
Analog & Physical modeling & Simulates analog circuits via real-time math; no wavetables; warm, vintage character \\
Collision & Physical modeling & Mallet percussion synthesis; two resonator bodies \\
Electric & Physical modeling & Electric piano tine/fork mechanism; modeled in real time \\
Tension & Physical modeling & String synthesis; bow, hammer, finger excitation models \\
Operator & FM synthesis & 4-operator FM with subtractive filter and additive mix \\
Wavetable & Wavetable & 2 oscillators + sub; wavetable morphing; analog-modeled filters \\
Drift & Subtractive / MPE & MPE-capable subtractive; per-note pressure and slide \\
Meld & Macro-oscillator / MPE & Bi-timbral; 2 macro oscillators with swappable engines; Mutable-inspired; 12-voice MPE \\
Granulator III & Granular & Robert Henke's device; real-time audio capture; MPE; glissando and vibrato \\
Sampler & Multi-sample & Multi-layer sampling; extensive modulation \\
Simpler & Single-sample & Slice, Classic, 1-Shot modes; granular pitch in 1-Shot \\
Impulse & Drum machine & 8-slot sample-based drum machine; per-slot tuning/stretch \\
Drum Rack & Drum framework & Maps MIDI notes to instrument chains; 128-note capacity \\
Drum Synths & Synthesis drums & 8 devices for kick, snare, hi-hat, etc.\ via synthesis \\
Bass & Monophonic analog & Virtual analog bass; monophonic only \\
Poli & Virtual analog & Subtractive + FM combination; polyphonic \\
CV Instrument & Hardware bridge & Sends CV/gate signals to modular/analog gear \\
CV Triggers & Hardware bridge & Gate triggers for modular/analog \\
\bottomrule
\end{tabularx}
\end{center}

\subsubsection{Live 12 New Instruments and Effects}

\textbf{Meld} (new in Live 12): A bi-timbral, 12-voice MPE-capable synthesizer with two
independent macro oscillators. Each oscillator offers multiple synthesis ``engines''---modes
with different underlying algorithms (analog, wavetable, spectral, FM-like). Two macro controls
per oscillator provide immediate sound-shaping without exposing underlying complexity. According
to CDM's pre-release coverage (November 2023), Meld was designed to bridge accessible
entry-level exploration with deep modular-inspired modulation---``for both ends of the
spectrum.'' Its design was influenced by Mutable Instruments' Eurorack modules.

\textbf{Roar} (new in Live 12): A multi-stage saturation and coloring effect. Three saturation
stages configurable in serial, parallel, mid/side, or multiband modes. Built-in compressor and
feedback routing. Full modulation matrix. The Ableton product page describes it as offering
``subtle warmth to wild and unpredictable sound degradation.''

\textbf{Granulator III} (updated in Live 12): Third generation of Robert Henke's granular
synthesis device. New MPE capability enables per-note bends, glissando, and vibrato.
Real-time audio capture added---start granularizing incoming audio immediately without
pre-recording.

\subsubsection{MPE Capability Summary}

\begin{center}
\rowcolors{2}{rowgray}{white}
\begin{tabularx}{\textwidth}{L{3cm}C{2.5cm}X}
\toprule
\rowcolor{primary}
\tableheader{Instrument} & \tableheader{MPE Capable} & \tableheader{Notes} \\
\midrule
Meld & Yes & 12-voice polyphonic; per-note pressure, slide, bend \\
Drift & Yes & Per-note modulation via MPE dimensions \\
Granulator III & Yes (v3) & Bend, glissando, vibrato per note \\
Wavetable & Partial & MPE support added in Live 11 \\
Analog & No & Physical modeling; polyphonic but not MPE-native \\
Operator & No & FM engine; polyphonic but not MPE-native \\
\bottomrule
\end{tabularx}
\end{center}

\subsection{Module 03 Reference: Hardware Ecosystem}

\subsubsection{Push 3 --- Two Configurations}

Push 3 (released May 2023) ships in two configurations:

\textbf{Push 3 Controller:} Functions as a USB-powered control surface for Live running on a
connected computer. Full feature parity with the standalone version for all controller
functions. Significantly lower price point. Recommended entry point according to trained
educators at Future Sound Academy.

\textbf{Push 3 Standalone:} Contains an Intel NUC Compute Element (processor + RAM + WiFi
in a credit-card-sized module). Runs a version of Live 12 Suite onboard without a computer.
Battery life approximately 2.5 hours. Built-in 2-channel audio interface with 6.35\,mm balanced
I/O. ADAT input/output sockets. Two stereo CV/gate signals (four mono). MIDI in/out. USB-A host
port for class-compliant MIDI devices. The compute element and battery are user-replaceable.

\textbf{Key Push 3 hardware features:}
\begin{itemize}[leftmargin=1.5em]
  \item 64 newly-developed MPE-capable pads (finger pressure and placement across X and Y axes)
  \item Full-color display strip across the top
  \item Large center encoder (new vs.\ Push 2)
  \item Wi-Fi for over-the-air software updates and Ableton Link
  \item Engineered originally in collaboration with Akai Professional (Push 1)
\end{itemize}

\textbf{Known limitation (as of early 2026):} Push 3 Standalone lacks a full Arrangement View,
no Follow Actions, and no per-Scene time signatures or tempos. Community response has included
a Max for Live--based arranger workaround (shared via Gumroad) that reconstructs basic
arrangement functionality in pure Max without JavaScript.

\subsubsection{Ableton Move (October 2024)}

A compact standalone groovebox announced October 8, 2024. Features include: battery power,
headphone socket, audio input, USB host port, USB-C. Positioned below Push 3 in the hardware
lineup. Designed as a companion device for sketching ideas that export directly into Ableton
Live. Community response noted limitations when used simultaneously with Push 3 (two devices
cannot independently control two different tracks).

\subsubsection{Ableton Link Protocol}

Ableton Link is a network synchronization technology integrated into Live and a growing
ecosystem of iOS and desktop applications. Key behaviors:
\begin{itemize}[leftmargin=1.5em]
  \item Keeps devices on the same local network in tempo and quantization phase
  \item Each device starts and stops independently (no master/slave)
  \item First app to join a session sets initial tempo; any participant can change tempo
  \item When multiple participants change tempo simultaneously, last-write wins
  \item MIDI Clock and MIDI Timecode also supported for non-Link hardware
\end{itemize}

\subsection{Module 04 Reference: Extensibility Layer}

\subsubsection{Max for Live (M4L)}

Max for Live integrates Cycling '74's Max/MSP visual programming environment directly into
Live's device chain. M4L devices have full access to the Live Object Model (LOM) and can
observe, query, and control any Live parameter. Key access mechanisms:

\begin{itemize}[leftmargin=1.5em]
  \item \texttt{live.object}, \texttt{live.observer}, \texttt{live.remote} Max objects
  \item JavaScript support via \texttt{js} object (often preferred for complex LOM traversal)
  \item Max for Live devices appear identically to native Live devices in the device chain
  \item Live 12 adds ``Edit in Max'' option from device title bar context menu
  \item Performance Pack (Live 12) ships four new M4L devices for live performance
  \item LFO, Shaper, Envelope Follower M4L devices now allow offset adjustment of modulated
  destinations---a significant usability improvement in Live 12
\end{itemize}

\subsubsection{Python MIDI Remote Scripts}

Live's MIDI Remote Scripts are Python modules that configure control surfaces (hardware
controllers). Key technical facts:
\begin{itemize}[leftmargin=1.5em]
  \item Live ships Python bytecode (\texttt{.pyc}) files for all supported controllers
  \item Community member Julien Bayle (structure-void.com) decompiled these scripts beginning
  in 2011, enabling unofficial modification and extension
  \item Live 11 upgraded from Python 2 to Python 3; all custom scripts must use Python 3
  \item Scripts live at: \texttt{Users/[user]/Music/Ableton/User Library/Remote Scripts} (macOS)
  \item Full unofficial API documentation maintained by NSUSpray and Julien Bayle
  \item Ableton's co-founder Robert Henke acknowledged the decompilation community early on:
  community reports quote him as saying ``we know you can. do what you can but we won't support it''
\end{itemize}

\subsubsection{AbletonOSC}

AbletonOSC (GitHub: ideoforms/AbletonOSC) is an open-source MIDI Remote Script providing
comprehensive OSC (Open Sound Control) access to the full LOM:

\begin{center}
\rowcolors{2}{rowgray}{white}
\begin{tabularx}{\textwidth}{L{3.5cm}X}
\toprule
\rowcolor{primary}
\tableheader{Object Type} & \tableheader{Key Capabilities via OSC} \\
\midrule
Song & Tempo, track names, scene list, transport, global properties \\
Track & Volume, pan, mute, solo, arm, send levels, clip slot queries \\
Clip & Launch/stop, note query/add/remove, playing position, quantization \\
Device & Parameter names and values, device enable/disable \\
MIDI CC & Create CC-to-parameter assignments \\
\bottomrule
\end{tabularx}
\end{center}

AbletonOSC listens on port 11000 by default. A command-line console utility (\texttt{run-console.py})
ships with the project for rapid testing.

\subsubsection{Extensibility Decision Tree}

\begin{center}
\rowcolors{2}{rowgray}{white}
\begin{tabularx}{\textwidth}{L{4.5cm}L{2.5cm}X}
\toprule
\rowcolor{primary}
\tableheader{Goal} & \tableheader{Recommended Path} & \tableheader{Rationale} \\
\midrule
Build a custom instrument or effect & Max for Live & Full audio/MIDI processing in device chain \\
Control Live parameters from external app & AbletonOSC & Network-accessible, no Live restart required \\
Deep controller customization & Python Scripts & Closest to hardware; handles LED feedback \\
No-code controller mapping & UserConfiguration.txt & Simplest path for basic transport/track control \\
Automate Live from command line & AbletonOSC & Platform-agnostic; works over localhost or LAN \\
\bottomrule
\end{tabularx}
\end{center}

\subsection{Module 05 Reference: Education, Community, and Ecosystem}

\subsubsection{Certified Trainer Program}

Launched by Ableton in 2008. As of the most recent program data: more than 360 certified
trainers operating across 50 countries in 38 languages. Certification requires an in-person
event at a location sourced by Ableton; online certification is explicitly not offered.
Candidates must demonstrate teaching and Live skills over several days of presentations.

\textbf{Notable training institutions:}
\begin{itemize}[leftmargin=1.5em]
  \item \textbf{343 Labs} (New York City and Berlin): Production and performance community with
  multiple Certified Trainers; free weekly masterclasses via Discord and YouTube
  \item \textbf{Point Blank} (London, Los Angeles, Mumbai, Ibiza, China, Online): Named
  ``Best Music Production School'' by DJ Mag; offers 2-year and 4-year degree programs
  alongside Live training
  \item \textbf{Delta Sound Labs} (Online, in collaboration with Case Western Reserve University,
  Cleveland): Intensive online production education since 2018; synchronous winter/summer sessions
\end{itemize}

\subsubsection{Free Learning Resources}

\begin{itemize}[leftmargin=1.5em]
  \item \textbf{Learning Music} (learningmusic.ableton.com): Browser-based synthesizer and
  interactive lessons covering fundamentals of music making; no software installation required
  \item \textbf{74 Creative Strategies for Electronic Producers}: Free online resource; tips
  from Ableton's creative blog
  \item \textbf{Ableton Reference Manual v12}: Comprehensive official documentation; among the
  most thorough DAW manuals in the industry
  \item \textbf{Loop Summit}: Ableton's annual conference for music makers; archived talks
  available online
\end{itemize}

\subsubsection{Historical Context}

\begin{itemize}[leftmargin=1.5em]
  \item Live created by Gerhard Behles, Robert Henke, and Bernd Roggendorf in the mid-1990s
  \item Ableton described as having ``made it easier for musicians to use computers as instruments
  in live performance without programming''---Wikipedia credits this with contributing to the
  rise of global festival culture in the 2000s
  \item 2010: Max for Live introduced, enabling direct Max/MSP integration
  \item 2016: Robert Henke left Ableton to focus on his artistic project Monolake; his
  Granulator III device continues to ship with Live Suite
  \item Push 1 (engineered by Akai Professional), Push 2 (2015), Push 3 (2023), Move (2024)
  mark a hardware evolution from controller to standalone instrument
\end{itemize}

\subsection{Safety and Sensitivity Considerations}

\begin{center}
\rowcolors{2}{rowgray}{white}
\begin{tabularx}{\textwidth}{L{3cm}L{2cm}X}
\toprule
\rowcolor{primary}
\tableheader{Area} & \tableheader{Level} & \tableheader{Guidance} \\
\midrule
Source attribution & GATE & All instrument specs cite Ableton Reference Manual v12 or official product pages \\
API documentation & ANNOTATE & Python decompilation community is unofficial; Ableton does not support it \\
Pricing data & ANNOTATE & Prices change; cite at time of research and advise verification \\
Push 3 Standalone limits & GATE & Document limitation status accurately per date; note community workarounds with appropriate caveats \\
Robert Henke attribution & STANDARD & Granulator III is his device; credit accurately \\
\bottomrule
\end{tabularx}
\end{center}

\subsection{Source Bibliography}

\subsubsection{Primary: Official Ableton Documentation}
\begin{itemize}[leftmargin=1.5em]
  \item Ableton AG. \textit{Ableton Live Reference Manual, Version 12}. ableton.com/en/manual/ (2024--2026)
  \item Ableton AG. \textit{Session View --- Reference Manual v12}. ableton.com/en/manual/session-view/
  \item Ableton AG. \textit{Arrangement View --- Reference Manual v12}. ableton.com/en/manual/arrangement-view/
  \item Ableton AG. \textit{Live Instrument Reference --- Manual v12}. ableton.com/en/manual/live-instrument-reference/
  \item Ableton AG. \textit{Synchronizing with Link, Tempo Follower, and MIDI --- Manual v12}
  \item Ableton AG. \textit{Live 12 Release Notes}. ableton.com/en/release-notes/live-12/
  \item Ableton AG. \textit{What's New in Live 12}. ableton.com/en/live/
  \item Ableton AG. \textit{Certification Program}. ableton.com/en/education/certification-program/
  \item Ableton AG. \textit{Push 3 Manual}. ableton.com/en/push/manual/
  \item Ableton AG. \textit{Meld: A Look at Live 12's New Bi-Timbral Synth} (blog post, 2024)
  \item Ableton AG. \textit{Installing Third-Party Remote Scripts}. help.ableton.com
  \item Ableton AG. \textit{Creating Your Own Control Surface Script}. help.ableton.com
  \item Ableton AG. \textit{Controlling Live Using Max for Live}. help.ableton.com
\end{itemize}

\subsubsection{Professional and Technical Publications}
\begin{itemize}[leftmargin=1.5em]
  \item Sherburne, Philip. ``Ableton Live 12: A Guide to Everything That's New.'' \textit{CDM Create Digital Music}, November 17, 2023. cdm.link
  \item Ableton Live. \textit{Wikipedia}. Accessed March 2026. (Historical facts and edition lists)
  \item \textit{Resident Advisor}. ``Ableton Reveals New Features in Live 12.'' ra.co/news/79804, November 2023
  \item \textit{Synth Anatomy}. ``Ableton Push 3 --- The Popular Live MIDI Controller Goes Standalone with MPE.'' May 23, 2023
  \item \textit{Sound on Sound}. ``Ableton Live: Session \& Arrangement Views.'' soundonsound.com
  \item Future Sound Academy. ``Ableton Push 3 Review 2026: Standalone vs Controller.'' 2026
\end{itemize}

\subsubsection{Open Source and Community}
\begin{itemize}[leftmargin=1.5em]
  \item Fisher, Ian (ideoforms). \textit{AbletonOSC}. GitHub: ideoforms/AbletonOSC (MIDI Remote Script for OSC control)
  \item Bayle, Julien. \textit{Ableton Live 11 MIDI Remote Scripts (Decompiled)}. GitHub: gluon/AbletonLive11\_MIDIRemoteScripts; structure-void.com
  \item NSUSpray. \textit{Unofficial Live API Documentation}. nsuspray.github.io/Live\_API\_Doc/
\end{itemize}

% =========================================================
% STAGE 3 — MISSION PACKAGE
% =========================================================
\newpage
\stageheader{STAGE 3 --- MISSION PACKAGE}{tertiary}

\section{Milestone Specification}

\subsection{Mission Objective}

Produce a comprehensive, self-contained research document on Ableton Live 12 covering its
core dual-view architecture, 20-instrument synthesis arsenal, hardware ecosystem (Push 3,
Move, Ableton Link), extensibility layer (M4L, Python scripts, AbletonOSC/LOM), and education
community. The deliverable serves as both a standalone reference and as the research foundation
for integrating Ableton Live into the GSD ecosystem's audio analysis, College of Knowledge,
and GSD-OS audio layer.

\subsection{Architecture Overview}

\begin{asciibox}
MISSION: ABLETON LIVE DEEP RESEARCH
=====================================

[WAVE 0: FOUNDATION]
  Haiku: Shared schemas, source index, bibliography template, test scaffolding
  Output: ableton_schema.json, source_index.md

[WAVE 1-A: CORE + INSTRUMENTS] (parallel)          [WAVE 1-B: HARDWARE + EXT] (parallel)
  Sonnet (EXEC-1 + CRAFT-ARCH):                      Sonnet (EXEC-2 + CRAFT-EXT):
    - Module 01: Core Architecture doc                  - Module 03: Hardware Ecosystem doc
    - Module 02: Synthesis Arsenal doc                  - Module 04: Extensibility Layer doc
  Opus: Domain analysis + LOM deep dive              Sonnet: LOM hierarchy mapping
  Output: mod01.md, mod02.md                         Output: mod03.md, mod04.md

[WAVE 1-C: EDUCATION] (parallel with 1-A and 1-B)
  Sonnet (EXEC): Module 05: Education + Community doc
  Output: mod05.md

[CAPCOM GATE]: Human review of five module docs -- approve before synthesis

[WAVE 2: SYNTHESIS]
  Opus (INTEG + FLIGHT):
    - Cross-module connections (instrument --> hardware MPE; LOM --> GSD DACP)
    - Synthesis taxonomy (architecture --> Space Between math)
    - GSD integration pathway document
  Output: synthesis_notes.md, integration_map.md

[WAVE 3: PUBLICATION + VERIFICATION]
  Sonnet (VERIFY + LOG):
    - Source validation against all module citations
    - Final compiled PDF (XeLaTeX, 3 passes)
    - Index HTML generation
  Output: ableton_live_mission_final.pdf, index.html

TOTAL: 3 parallel track waves + 1 synthesis + 1 publication
\end{asciibox}

\subsection{System Layers}

\begin{enumerate}
  \item \textbf{Foundation Layer}: Shared schema, source index, test scaffolding (Wave 0, Haiku)
  \item \textbf{Research Layer A}: Core architecture and synthesis arsenal (Wave 1-A, Sonnet/Opus)
  \item \textbf{Research Layer B}: Hardware ecosystem and extensibility (Wave 1-B, Sonnet)
  \item \textbf{Research Layer C}: Education and community (Wave 1-C, Sonnet)
  \item \textbf{Synthesis Layer}: Cross-module integration and GSD connection analysis (Wave 2, Opus)
  \item \textbf{Publication Layer}: Compiled document and verification (Wave 3, Sonnet)
\end{enumerate}

\subsection{Deliverables}

\begin{center}
\rowcolors{2}{rowgray}{white}
\begin{tabularx}{\textwidth}{C{0.5cm}L{4cm}L{3.5cm}X}
\toprule
\rowcolor{primary}
\tableheader{\#} & \tableheader{Deliverable} & \tableheader{Acceptance Criteria} & \tableheader{Spec} \\
\midrule
1 & Module 01: Core Architecture & Session/Arrangement model fully documented & Markdown, sourced to Ableton Manual v12 \\
2 & Module 02: Synthesis Arsenal & All 20 instruments catalogued with taxonomy & Table + prose, sourced to Ableton Manual v12 \\
3 & Module 03: Hardware Ecosystem & Push 3, Move, Link documented with limitations & Markdown, sourced to Ableton product pages \\
4 & Module 04: Extensibility Layer & M4L/Python/OSC decision tree produced & Markdown with decision table \\
5 & Module 05: Education Community & Trainer program, institutions, free resources & Markdown, sourced to ableton.com \\
6 & Synthesis Document & Cross-module connections + GSD pathways & Markdown, Opus-produced \\
7 & Final PDF & All 5 modules + synthesis compiled, LaTeX & XeLaTeX, 3-pass compilation \\
8 & Index HTML & Download portal with usage instructions & HTML with color scheme matching PDF \\
\bottomrule
\end{tabularx}
\end{center}

\subsection{Component Breakdown}

\begin{center}
\rowcolors{2}{rowgray}{white}
\begin{tabularx}{\textwidth}{L{3.5cm}L{2.5cm}C{2cm}C{2cm}}
\toprule
\rowcolor{primary}
\tableheader{Component} & \tableheader{Dependencies} & \tableheader{Model} & \tableheader{Est.\ Tokens} \\
\midrule
Schema + source index & None & Haiku & 2K \\
Module 01 (Core Arch) & Schema & Sonnet & 12K \\
Module 02 (Synthesis) & Schema & Sonnet & 15K \\
Module 03 (Hardware) & Schema & Sonnet & 10K \\
Module 04 (Extension) & Schema & Sonnet & 12K \\
Module 05 (Education) & Schema & Sonnet & 8K \\
LOM deep analysis & Mod 04 & Opus & 10K \\
Synthesis doc & Mod 01--05 & Opus & 15K \\
GSD integration map & Synthesis & Opus & 8K \\
PDF compilation & All modules & Sonnet & 6K \\
Index HTML & PDF & Sonnet & 3K \\
\midrule
\multicolumn{3}{l}{\textbf{Total Estimated}} & \textbf{\textasciitilde 101K} \\
\bottomrule
\end{tabularx}
\end{center}

\subsection{Model Assignment Rationale}

\begin{center}
\rowcolors{2}{rowgray}{white}
\begin{tabularx}{\textwidth}{L{2cm}L{2cm}C{2.5cm}X}
\toprule
\rowcolor{primary}
\tableheader{Role} & \tableheader{Agent} & \tableheader{Model} & \tableheader{Rationale} \\
\midrule
FLIGHT & Orchestrator & Opus (\textasciitilde 30\%) & Mission direction; cross-module synthesis judgment; GSD connection analysis \\
EXEC-1 & Document prod & Sonnet (\textasciitilde 60\%) & Structured content assembly; Module 01--02 \\
EXEC-2 & Document prod & Sonnet & Structured content assembly; Module 03--04 \\
EXEC-3 & Document prod & Sonnet & Module 05 + Index HTML \\
CRAFT-ARCH & Domain spec & Opus & Core architecture and LOM deep analysis \\
CRAFT-EXT & Domain spec & Sonnet & Extensibility layer technical mapping \\
SCOUT & Research & Haiku (\textasciitilde 10\%) & Web search gap-fill; source validation \\
VERIFY & Verification & Sonnet & Citation tracing; accuracy checking \\
INTEG & Integration & Opus & Cross-module relationship validation \\
BUDGET & Resources & Haiku & Token tracking \\
SURGEON & Health & Sonnet & Research quality monitoring \\
LOG & Chronicle & Haiku & Commit messages; audit trail \\
\bottomrule
\end{tabularx}
\end{center}

\section{Wave Execution Plan}

\subsection{Wave Summary}

\begin{center}
\rowcolors{2}{rowgray}{white}
\begin{tabularx}{\textwidth}{C{1.5cm}L{3cm}C{2cm}C{2cm}X}
\toprule
\rowcolor{primary}
\tableheader{Wave} & \tableheader{Name} & \tableheader{Mode} & \tableheader{Model} & \tableheader{Outputs} \\
\midrule
0 & Foundation & Sequential & Haiku & Schema, source index, test scaffolding \\
1-A & Core + Instruments & Parallel & Sonnet/Opus & Module 01, Module 02, LOM analysis \\
1-B & Hardware + Extension & Parallel & Sonnet & Module 03, Module 04, decision tree \\
1-C & Education & Parallel & Sonnet & Module 05 \\
GATE & CAPCOM Review & --- & Human & 5-module review + approval \\
2 & Synthesis & Sequential & Opus & Synthesis doc, GSD integration map \\
3 & Publication & Sequential & Sonnet & Final PDF, Index HTML \\
\bottomrule
\end{tabularx}
\end{center}

\subsection{Cache Optimization Strategy}

\begin{itemize}[leftmargin=1.5em]
  \item System prompt (schema, source index, style guide) computed once in Wave 0 and
  reused across all Wave 1 agents---\textbf{5-minute TTL cache window} exploited by
  starting all three Wave 1 tracks within the same session
  \item Shared attribute layer (instrument list, edition tiers, edition-instrument mapping)
  produced in Wave 0 and injected into Modules 01 and 02 without re-derivation
  \item LOM hierarchy analysis (Wave 1-A Opus) cached and passed directly into Wave 2
  Synthesis agent---no re-traversal needed
  \item Module markdown files passed to Wave 2 as document context, not re-generated
\end{itemize}

\subsection{Token Budget}

\begin{center}
\rowcolors{2}{rowgray}{white}
\begin{tabularx}{\textwidth}{L{3cm}C{2cm}C{2cm}C{2cm}X}
\toprule
\rowcolor{primary}
\tableheader{Wave} & \tableheader{Input} & \tableheader{Output} & \tableheader{Total} & \tableheader{Notes} \\
\midrule
Wave 0 & 5K & 5K & 10K & Foundation only \\
Wave 1-A & 20K & 30K & 50K & Core + Instruments, LOM \\
Wave 1-B & 20K & 25K & 45K & Hardware + Extension \\
Wave 1-C & 10K & 10K & 20K & Education \\
Wave 2 & 50K & 25K & 75K & Synthesis (large context) \\
Wave 3 & 30K & 10K & 40K & Publication + verification \\
\midrule
\multicolumn{3}{l}{\textbf{Grand Total}} & \textbf{240K} & \textbf{Across 3 sessions} \\
\bottomrule
\end{tabularx}
\end{center}

\subsection{Dependency Graph}

\begin{asciibox}
[Wave 0: Schema] ──────────────────────────────────────────────┐
        |                                                       |
        ├──────────── [Wave 1-A] ─────────────────────────┐    |
        |   Mod01 (Sonnet) ──> LOM Analysis (Opus)        |    |
        |   Mod02 (Sonnet)                                 |    |
        |                                                  |    |
        ├──────────── [Wave 1-B] ─────────────────────┐   |    |
        |   Mod03 (Sonnet)                            |   |    |
        |   Mod04 (Sonnet) ──> Decision Tree          |   |    |
        |                                             |   |    |
        └──────────── [Wave 1-C] ─────────────────┐  |   |    |
                     Mod05 (Sonnet)                |  |   |    |
                                                   |  |   |    |
                           [CAPCOM GATE]           |  |   |    |
                               |                   |  |   |    |
                    [Wave 2: Synthesis] <───────────┘──┘───┘    |
                         Opus: Cross-module                      |
                         Opus: GSD integration map               |
                               |                                 |
                    [Wave 3: Publication] <───────────────────────┘
                         Sonnet: XeLaTeX PDF
                         Sonnet: Index HTML
\end{asciibox}

\section{Test Plan}

\subsection{Test Categories}

\begin{center}
\rowcolors{2}{rowgray}{white}
\begin{tabularx}{\textwidth}{L{3.5cm}XC{1.5cm}C{1.5cm}}
\toprule
\rowcolor{primary}
\tableheader{Category} & \tableheader{Purpose} & \tableheader{Count} & \tableheader{Action} \\
\midrule
Safety-critical & Non-negotiable source and attribution boundaries & 4 & BLOCK \\
Core functionality & Module deliverable acceptance criteria & 20 & BLOCK \\
Integration & Cross-module and GSD-connection validation & 6 & BLOCK \\
Edge cases & Robustness and currency & 4 & LOG \\
\midrule
\multicolumn{2}{l}{\textbf{Total}} & \textbf{34} & \\
\bottomrule
\end{tabularx}
\end{center}

\subsection{Safety-Critical Tests}

\begin{center}
\rowcolors{2}{rowgray}{white}
\begin{tabularx}{\textwidth}{C{1.5cm}L{3.5cm}X}
\toprule
\rowcolor{primary}
\tableheader{ID} & \tableheader{Verifies} & \tableheader{Expected Behavior} \\
\midrule
SC-SRC & Source quality & All instrument specs, hardware specs, and API documentation claims trace to Ableton's official manual, official product pages, CDM, or open-source GitHub repositories. Zero entertainment media. \\
SC-NUM & Numerical attribution & Edition counts (20 instruments, 58 audio effects, 14 MIDI effects, 71\,GB content, 360+ trainers, 50 countries, 38 languages) each attributed to specific source \\
SC-ADV & No advocacy & Document presents Ableton Live's features factually; no recommendation over competing DAWs; push-back on limitations documented objectively \\
SC-API & API unofficial status & Python MIDI Remote Script decompilation and AbletonOSC clearly noted as unofficial/community tools; no claim that Ableton supports them \\
\bottomrule
\end{tabularx}
\end{center}

\subsection{Core Functionality Tests (Selected)}

\begin{center}
\rowcolors{2}{rowgray}{white}
\begin{tabularx}{\textwidth}{C{1.5cm}L{4cm}X}
\toprule
\rowcolor{primary}
\tableheader{ID} & \tableheader{Verifies} & \tableheader{Expected Behavior} \\
\midrule
CF-01 & Session/Arrangement exclusivity & Document explicitly states clip exclusivity rule and ``Back to Arrangement'' state transition \\
CF-02 & Warping engine & Six Warp Modes documented with use-cases \\
CF-03 & Instrument taxonomy completeness & All 20 Suite instruments listed with synthesis architecture \\
CF-04 & MPE instrument identification & Table identifies which instruments support MPE and at what level \\
CF-05 & Push 3 configuration accuracy & Standalone vs.\ Controller differences documented; Intel NUC cited \\
CF-06 & Push 3 limitation currency & Arrangement View absence noted with date qualifier \\
CF-07 & Ableton Move documented & October 2024 release; key form factor and capabilities \\
CF-08 & Link protocol behavior & Tempo negotiation, start/stop independence, quantization documented \\
CF-09 & M4L LOM access documented & At least three LOM object types with key properties listed \\
CF-10 & Extensibility decision tree & Table or flowchart mapping integration goal to recommended mechanism \\
CF-11 & Trainer program scale & 360+, 50 countries, 38 languages, since 2008 --- all four figures cited \\
CF-12 & Free resources listed & Learning Music browser app and Reference Manual both identified \\
\bottomrule
\end{tabularx}
\end{center}

\subsection{Integration Tests}

\begin{center}
\rowcolors{2}{rowgray}{white}
\begin{tabularx}{\textwidth}{C{1.5cm}L{4cm}X}
\toprule
\rowcolor{primary}
\tableheader{ID} & \tableheader{Verifies} & \tableheader{Expected Behavior} \\
\midrule
IN-01 & MPE hardware--instrument link & Push 3 MPE pads explicitly connected to MPE-capable instruments (Meld, Drift, Granulator III) \\
IN-02 & Extensibility--hardware link & AbletonOSC access documented as applicable to Push 3 Standalone via network \\
IN-03 & GSD subagent analogy & Synthesis document draws explicit connection between Session View clip model and GSD subagent architecture \\
IN-04 & Space Between math link & Operator FM synthesis linked to complex number modulation from \textit{The Space Between} \\
IN-05 & Cross-module data consistency & Instrument counts in Module 02 match edition table in Module 01 \\
IN-06 & Self-containment & A reader with no prior Ableton knowledge can understand the full ecosystem from Stage 1 alone \\
\bottomrule
\end{tabularx}
\end{center}

\subsection{Verification Matrix}

\begin{center}
\rowcolors{2}{rowgray}{white}
\begin{tabularx}{\textwidth}{XL{3.5cm}C{1.5cm}}
\toprule
\rowcolor{primary}
\tableheader{Success Criterion} & \tableheader{Test ID(s)} & \tableheader{Status} \\
\midrule
1. Session/Arrangement model fully documented & CF-01, IN-06 & Pending \\
2. All 20 instruments catalogued & CF-03, CF-04, IN-05 & Pending \\
3. Extensibility decision tree produced & CF-09, CF-10 & Pending \\
4. LOM documented to three levels & CF-09, IN-02 & Pending \\
5. Push 3 limitations documented & CF-05, CF-06 & Pending \\
6. Link protocol documented & CF-08 & Pending \\
7. Education ecosystem mapped & CF-11, CF-12 & Pending \\
8. GSD integration pathways & IN-03, IN-04 & Pending \\
9. All claims sourced to official docs & SC-SRC, SC-NUM & Pending \\
10. Document self-contained & IN-06 & Pending \\
11. Synthesis taxonomy with math link & CF-03, IN-04 & Pending \\
12. Test plan covers all modules & All CF, IN & Pending \\
\bottomrule
\end{tabularx}
\end{center}

\section{Mission Crew Manifest}

\begin{center}
\rowcolors{2}{rowgray}{white}
\begin{tabularx}{\textwidth}{L{2cm}L{3cm}X}
\toprule
\rowcolor{primary}
\tableheader{Role} & \tableheader{Agent Name} & \tableheader{Primary Responsibility} \\
\midrule
FLIGHT & Orchestrator & Mission direction; go/no-go gates; GSD synthesis \\
PLAN & Planner & Wave decomposition; dependency management; cache planning \\
SCOUT & Research Agent & Source gathering; gap-fill web research; bibliography \\
EXEC-1 & Sonnet-Alpha & Module 01 + 02 document production \\
EXEC-2 & Sonnet-Beta & Module 03 + 04 document production \\
EXEC-3 & Sonnet-Gamma & Module 05 + Index HTML production \\
CRAFT-ARCH & Tahoma & Core architecture and Live Object Model deep analysis \\
CRAFT-EXT & Raven & Extensibility layer technical mapping \\
VERIFY & Verification & Source validation; accuracy checking against Ableton Manual \\
INTEG & Integration & Cross-module relationships; GSD pathway validation \\
CAPCOM & HITL Interface & Human approval at module-review gate \\
BUDGET & Resource Manager & Token budget tracking across 3 sessions \\
SURGEON & Health Monitor & Research quality and citation completeness monitoring \\
LOG & Osprey & Commit messages; milestone summaries; audit trail \\
\bottomrule
\end{tabularx}
\end{center}

\section{Execution Summary}

\begin{center}
\rowcolors{2}{rowgray}{white}
\begin{tabularx}{\textwidth}{L{5cm}X}
\toprule
\rowcolor{primary}
\tableheader{Metric} & \tableheader{Value} \\
\midrule
Mission name & Ableton Live: Architecture, Extensibility \& Ecosystem \\
Activation profile & Squadron (12 active roles; 14 crew total) \\
Research modules & 5 (Core Architecture, Instruments, Hardware, Extension, Education) \\
Wave count & 7 (0, 1-A, 1-B, 1-C, GATE, 2, 3) \\
Parallel tracks & 3 (Waves 1-A, 1-B, 1-C run simultaneously) \\
HITL gates & 1 (post Wave 1, pre Wave 2) \\
Estimated token budget & 240K across 3 sessions \\
Primary model & Sonnet (\textasciitilde 60\%) \\
Judgment model & Opus (\textasciitilde 30\%) \\
Scaffold model & Haiku (\textasciitilde 10\%) \\
Deliverable count & 8 (5 modules + synthesis + PDF + HTML) \\
Handoff target & gsd-skill-creator dev branch; \texttt{.planning/staging/inbox/} \\
GSD milestone context & Post-v1.50 audio integration layer \\
Safety tests & 4 mandatory-pass / BLOCK \\
Core tests & 20 mandatory-pass / BLOCK \\
Integration tests & 6 mandatory-pass / BLOCK \\
Edge case tests & 4 best-effort / LOG \\
\bottomrule
\end{tabularx}
\end{center}

\vspace{2em}

\begin{quotebox}
\textit{``In Live's Arrangement View, as in all traditional sequencing programs, everything
happens along a fixed song timeline. For a number of applications, this is a limiting paradigm.
This is exactly what Live's unique Session View is for.''}

\vspace{0.5em}
{\small\sffamily\textcolor{subtlegray}{--- Ableton Reference Manual, Version 12 (paraphrased)}}

\vspace{1em}
\noindent The space between the Session's open field and the Arrangement's committed timeline
is not a limitation to be engineered away. It is the creative gap where music---and
agents---find their most interesting decisions. In the GSD ecosystem, every subagent lives
in this space: bounded by its context window (a clip), free to launch non-linearly (Session
View), and eventually committed to the pipeline (Arrangement View). Ableton Live did not
invent this architecture. It discovered it.
\end{quotebox}

\end{document}
